\documentclass[12pt]{article}
\usepackage[utf8]{inputenc}
\usepackage[T1]{fontenc}
\usepackage[margin=1in]{geometry}
\usepackage{setspace}
\usepackage{xcolor}
\usepackage[colorlinks=true, urlcolor=blue!60!black, linkcolor=blue!60!black]{hyperref}
\setlength{\parindent}{0em}
\setlength{\parskip}{0.45em}
\renewcommand{\baselinestretch}{1.0}
\begin{document}
\begin{center}
{\Large \textit{Rust v. Sullivan}, 500 U.S. 173 (1991)\par}
\vspace{0.3em}
{\normalsize Case Brief --- One-Page Edition\par}
\end{center}

\textbf{Memory Jogger.} The Title X gag rule: government may define the limits of the programs it funds (\href{https://supreme.justia.com/cases/federal/us/500/173/}{500 U.S. 173}).

\textbf{Facts.} Title X funds family planning but not ``programs where abortion is a method of family planning.'' In 1988 HHS barred Title X projects from abortion counseling, referral, and advocacy and required physical and financial separation from abortion activities. Grantees and physicians sued.

\textbf{Procedural History.} The Second Circuit upheld the rules; the First and Tenth Circuits had struck them down.

\textbf{Votes.} Affirmed, 5--4. Rehnquist wrote for the Court; Blackmun, Stevens, and O'Connor dissented.

\textbf{Holding.} The rules are a permissible reading of an ambiguous statute and do not violate the First or Fifth Amendments: the government may fund one activity to the exclusion of another, limiting the funded project but not the grantee's speech outside it.

\textbf{Reasoning.} Under \textit{Maher} and \textit{Harris v. McRae}, selective funding is not viewpoint discrimination; grantees could still perform and advocate abortion through separate programs. Blackmun saw viewpoint-based intrusion into the doctor--patient relationship; O'Connor would have avoided the constitutional question.

\textbf{Examples.} \emph{Same-Side:} a smoking-cessation grant barring advice on vaping. \emph{Opposite-Side:} an organization-wide policy pledge required of grantees (\textit{AID v. AOSI}). \emph{Fence-Sitter:} speech limits on university research grants.

\textbf{Critique.} Defenders cite democratic control of public money; critics say the decision lets government buy silence and co-opt professionals.

\textbf{Key Quotation.} ``The Government can, without violating the Constitution, selectively fund a program to encourage certain activities it believes to be in the public interest, without at the same time funding an alternate program which seeks to deal with the problem in another way. In so doing, the Government has not discriminated on the basis of viewpoint; it has merely chosen to fund one activity to the exclusion of the other.'' --- Chief Justice Rehnquist, Opinion of the Court, 500 U.S. at 193.

\end{document}
